\scanpage{97}

Soit $G$ un schéma en gr. affine sur corps $k$ $j:G_{m_k}-G$ un hom, i.e. une structure de $\otimes$-graduation sur le foncteur oubli $T$ sur $\operatorname{Rep}(G)$.

Cette $\otimes$-graduation définit une $\deleted{\unsure{\oplus}}\otimes$-filtration, et on peut regarder les ss-groupes

\[
U\subset P\deleted{\unsure{H}}\subset G\qquad
\deleted{\unsure{R}}P=\underline{\mathrm{Aut}}_{\otimes,\mathrm{filtr}\deleted{\unsure{,i}}}\,T,\quad
U=\underline{\mathrm{Aut}}_{\otimes,\mathrm{filt},\mathrm{id.\,sur\,Gr}}(T)
\]

Lorsque $G$ \inserted{est réductif sur $k$ de car $0$, et que} admet un module fidèle $V$, \deleted{\unsure{alors}} alors $V$ est gradué par $\mathrm{Gr}_{\deleted{\unsure{g}}j}^i(V)$, filtré par

\[
\mathrm{Filt}^{p}(V)=\sum_{i\geq p}\mathrm{Gr}_j^i(V)
\]

et $P\deleted{\unsure{H}}$ est le ss-gpe de $\underline{\mathrm{Aut}}_k(V)$ laissant \unsure{invariant|invariante} cette filtration, et $U$ le sous-gpe noyau de $P\longrightarrow\prod_i\underline{\mathrm{Aut}}_k\mathrm{Gr}_{\underline{j}}^i(V)$.

Dans le cas général, $P$ et $U$ sont des intersections de tels groupes $P_V,U_V$ ($V\in\mathrm{Ob}\,\operatorname{Rep}(\unsure{G|Q})$) ; si $P$ est de t.f. on peut prendre \uline{un} $V$ tel que $P_V=P,\ U_V=U$.

\deleted{\unsure{Prenons|Précisons}} Soit $\mathfrak g=\mathrm{Lie}\,G$, $\unsure{\xi|\zeta}\in\mathfrak g$ \deleted{\unsure{défini|définit}} l'image du générateur can de $\mathfrak t=\mathrm{Lie}\,G_m$ par $\mathrm{Lie}\,j:\deleted{\mathrm{Lie}\,G_m\to\mathfrak g}\inserted{\mathfrak t\longrightarrow\mathfrak g}$, $\mathfrak g=\sum_{\alpha\in\mathbb Z}\mathfrak g^\alpha$ la décomposition de $\mathfrak g$ suivant les car de $G_m$ opérant sur $\mathfrak g$, $V=\sum_{\beta\in\mathbb Z}V^\beta$ de $\overline{\mathrm m}$.

On a donc

\[
[\mathfrak g^\alpha,\mathfrak g^\beta]\subset\mathfrak g^{\alpha+\beta}
\qquad \mathfrak g^\alpha\cdot V^\beta\subset V^{\alpha+\beta}
\]

L'Algèbre de Lie $\mathfrak p$ de $P_V$ est formée des
% archive sheet number: 96
